% ============================================================
% 高考数学满分突破系列 · 第2讲
% 导数的深层理解
% 编译命令：xelatex -interaction=nonstopmode 高考满分突破2_导数的深层理解.tex
% ============================================================
\documentclass[12pt, a4paper]{ctexart}
\usepackage[margin=2.2cm]{geometry}
\usepackage{amsmath, amssymb}
\usepackage{enumitem}
\usepackage{tcolorbox}
\usepackage{booktabs}
\usepackage{array}
\usepackage{fancyhdr}
\usepackage{xcolor}
\usepackage{tikz}
\usepackage{hyperref}
\usepackage{graphicx}
\usepackage{multirow}
\usepackage{pgfplots}
\pgfplotsset{compat=1.18}

% 颜色定义
\definecolor{mainblue}{RGB}{41, 98, 168}
\definecolor{lightblue}{RGB}{230, 240, 255}
\definecolor{tipgreen}{RGB}{34, 139, 87}
\definecolor{lightgreen}{RGB}{235, 250, 240}
\definecolor{warnred}{RGB}{200, 50, 50}
\definecolor{lightred}{RGB}{255, 240, 240}
\definecolor{examorange}{RGB}{210, 120, 20}
\definecolor{lightorange}{RGB}{255, 248, 235}
\definecolor{apppurple}{RGB}{120, 60, 170}
\definecolor{lightpurple}{RGB}{245, 238, 255}

% 页眉页脚
\pagestyle{fancy}
\fancyhf{}
\fancyhead[L]{\small\textcolor{mainblue}{\textbf{高考数学满分突破 · 导数的深层理解}}}
\fancyhead[R]{\small\textcolor{mainblue}{第2讲}}
\fancyfoot[C]{\small —\ \thepage\ —}
\renewcommand{\headrulewidth}{0.6pt}
\renewcommand{\headrule}{\hbox to\headwidth{\color{mainblue}\leaders\hrule height \headrulewidth\hfill}}
\renewcommand{\footrulewidth}{0pt}

% tcolorbox 样式
\newtcolorbox{knowledgebox}[1][]{%
  colback=lightblue, colframe=mainblue,
  fonttitle=\bfseries\large, title=#1,
  arc=3mm, boxrule=0.8pt,
  left=8pt, right=8pt, top=6pt, bottom=6pt
}

\newtcolorbox{tipbox}[1][]{%
  colback=lightgreen, colframe=tipgreen,
  fonttitle=\bfseries, title=#1,
  arc=2mm, boxrule=0.6pt,
  left=6pt, right=6pt, top=4pt, bottom=4pt
}

\newtcolorbox{warnbox}[1][]{%
  colback=lightred, colframe=warnred,
  fonttitle=\bfseries, title=#1,
  arc=2mm, boxrule=0.6pt,
  left=6pt, right=6pt, top=4pt, bottom=4pt
}

\newtcolorbox{examplebox}[1][]{%
  colback=lightorange, colframe=examorange,
  fonttitle=\bfseries, title=#1,
  arc=3mm, boxrule=0.8pt,
  left=8pt, right=8pt, top=6pt, bottom=6pt
}

\newtcolorbox{appbox}[1][]{%
  colback=lightpurple, colframe=apppurple,
  fonttitle=\bfseries, title=#1,
  arc=3mm, boxrule=0.8pt,
  left=8pt, right=8pt, top=6pt, bottom=6pt
}

\setCJKmainfont{Hiragino Sans GB}
\setCJKsansfont{Heiti SC}

\begin{document}

% ============ 标题页 ============
\begin{center}
{\Huge\bfseries\textcolor{mainblue}{导数的深层理解}}\\[12pt]
{\Large\textcolor{gray}{高考数学满分突破\quad 第2讲}}\\[8pt]
\rule{0.6\textwidth}{1pt}\\[8pt]
{\normalsize 适用：高三冲刺\quad |\quad 课时：3课时}
\end{center}
\vspace{10pt}

% ============ 第一部分：知识点讲解 ============
\section*{\textcolor{mainblue}{一、知识点讲解}}

\subsection*{2.1\quad 局部线性化}

\begin{knowledgebox}[核心思想]
函数 $f(x)$ 在 $x = a$ 附近的最佳线性近似为：
\[
f(x) \approx f(a) + f'(a)(x - a)
\]
这就是导数的几何意义——用切线代替曲线。在 $x$ 越接近 $a$ 时，近似越精确。
\end{knowledgebox}

局部线性化的本质是"以直代曲"。在 $x = a$ 的邻域内：
\[
f(x) = f(a) + f'(a)(x-a) + o(x-a)
\]
其中 $o(x-a)$ 表示比 $(x-a)$ 更高阶的无穷小量。

\begin{tipbox}[高考应用]
局部线性化是理解导数概念的基础。高考中"用导数研究函数性质"的题目，本质上都是在利用这个思想。理解了局部线性化，就理解了为什么导数能判断单调性、求极值。
\end{tipbox}

\subsection*{2.2\quad 泰勒展开}

\begin{knowledgebox}[常用函数的泰勒展开（麦克劳林级数，即 $x_0 = 0$ 处展开）]
\begin{alignat*}{2}
e^x &= 1 + x + \frac{x^2}{2!} + \frac{x^3}{3!} + \cdots &&= \sum_{n=0}^{\infty} \frac{x^n}{n!}, \quad x \in \mathbb{R} \\[6pt]
\ln(1+x) &= x - \frac{x^2}{2} + \frac{x^3}{3} - \frac{x^4}{4} + \cdots &&= \sum_{n=1}^{\infty} \frac{(-1)^{n+1} x^n}{n}, \quad x \in (-1, 1] \\[6pt]
\sin x &= x - \frac{x^3}{3!} + \frac{x^5}{5!} - \cdots &&= \sum_{n=0}^{\infty} \frac{(-1)^n x^{2n+1}}{(2n+1)!}, \quad x \in \mathbb{R} \\[6pt]
\cos x &= 1 - \frac{x^2}{2!} + \frac{x^4}{4!} - \cdots &&= \sum_{n=0}^{\infty} \frac{(-1)^n x^{2n}}{(2n)!}, \quad x \in \mathbb{R} \\[6pt]
\frac{1}{1-x} &= 1 + x + x^2 + x^3 + \cdots &&= \sum_{n=0}^{\infty} x^n, \quad |x| < 1
\end{alignat*}
\end{knowledgebox}

\begin{warnbox}[重要提醒]
高考不要求写泰勒级数，但理解泰勒展开的前几项，可以帮助快速判断函数在某点附近的行为，从而秒杀许多放缩题和切线问题。\textbf{这是一把隐藏的武器}。
\end{warnbox}

\subsection*{2.3\quad 泰勒秒杀切线放缩}

\begin{knowledgebox}[常用放缩不等式（源自泰勒展开）]
\begin{enumerate}[leftmargin=2em]
  \item $e^x \geq 1 + x$，等号当且仅当 $x = 0$（由 $e^x$ 展开的前两项）
  \item $e^x \geq 1 + x + \dfrac{x^2}{2}$，等号当且仅当 $x = 0$（取前三项）
  \item $\ln(1+x) \leq x$，等号当且仅当 $x = 0$（即 $\ln t \leq t - 1$）
  \item $\ln(1+x) \geq x - \dfrac{x^2}{2}$，$x \geq 0$ 时成立
  \item $\sin x \leq x$，$x \geq 0$ 时成立
  \item $\sin x \geq x - \dfrac{x^3}{6}$，$x \geq 0$ 时成立
  \item $\cos x \leq 1 - \dfrac{x^2}{2} + \dfrac{x^4}{24}$（交替级数截断估计）
\end{enumerate}
\end{knowledgebox}

\begin{tipbox}[秒杀心法]
放缩的本质：泰勒展开取前 $n$ 项作为近似，奇数次截断和偶数次截断的放缩方向不同。记住这组不等式，高考压轴放缩题可以"秒杀"。
\end{tipbox}

\subsection*{2.4\quad 高考压轴题的展开视角}

\begin{knowledgebox}[泰勒视角解题三步法]
\begin{enumerate}[leftmargin=2em]
  \item \textbf{识别结构}：观察目标不等式中出现的函数组合，如 $e^x - 1 - x$、$\ln(1+x) - x$ 等。
  \item \textbf{展开定位}：将函数在关键点（通常是 $x = 0$ 或 $x = 1$）展开到需要的阶数。
  \item \textbf{严格证明}：用泰勒展开得到放缩方向后，再用导数方法给出严格证明（高考答卷需要）。
\end{enumerate}
\end{knowledgebox}

\begin{warnbox}[考场策略]
泰勒展开是\textbf{思考工具}，不是\textbf{书写工具}。在考场上，你可以用泰勒展开快速找到答案和证明方向，但书写时必须用导数方法严格证明。
\end{warnbox}

% ============ 第二部分：例题精讲 ============
\section*{\textcolor{examorange}{二、例题精讲}}

\begin{examplebox}[【例题1】泰勒展开求近似值]
利用泰勒展开的前三项，计算 $e^{0.1}$ 的近似值，并估计误差。
\end{examplebox}

\textbf{解：}$e^x \approx 1 + x + \dfrac{x^2}{2}$，代入 $x = 0.1$：
\[
e^{0.1} \approx 1 + 0.1 + \frac{0.01}{2} = 1.105
\]
误差由第三项控制：$|R| < \dfrac{|x|^3}{3!} \cdot e^{|x|} < \dfrac{0.001}{6} \times e^{0.1} < 0.0002$。

精确值 $e^{0.1} \approx 1.10517$，近似值 $1.105$ 的误差在万分之二以内。

\begin{examplebox}[【例题2】泰勒展开求近似值]
利用泰勒展开计算 $\ln 1.2$ 的近似值（取前三项）。
\end{examplebox}

\textbf{解：}$\ln(1+x) \approx x - \dfrac{x^2}{2} + \dfrac{x^3}{3}$，代入 $x = 0.2$：
\[
\ln 1.2 \approx 0.2 - \frac{0.04}{2} + \frac{0.008}{3} = 0.2 - 0.02 + 0.00267 = 0.18267
\]
精确值 $\ln 1.2 \approx 0.18232$，近似精度很高。

\begin{examplebox}[【例题3】放缩证明（2020全国I卷压轴改编）]
证明：当 $x > 0$ 时，$e^x > 1 + x + \dfrac{x^2}{2}$。
\end{examplebox}

\textbf{泰勒视角：}$e^x = 1 + x + \dfrac{x^2}{2} + \dfrac{x^3}{6} + \cdots$，后面全是正项，所以 $e^x > 1 + x + \dfrac{x^2}{2}$。

\textbf{严格证明：}设 $f(x) = e^x - 1 - x - \dfrac{x^2}{2}$，则
\[
f'(x) = e^x - 1 - x, \quad f''(x) = e^x - 1
\]
当 $x > 0$ 时，$f''(x) = e^x - 1 > 0$，故 $f'(x)$ 在 $(0, +\infty)$ 上递增。又 $f'(0) = 0$，故 $f'(x) > 0$。

因此 $f(x)$ 在 $(0, +\infty)$ 上递增，$f(0) = 0$，故 $f(x) > 0$，即 $e^x > 1 + x + \dfrac{x^2}{2}$。\qed

\begin{examplebox}[【例题4】放缩证明]
证明：当 $x > 0$ 时，$\ln(1+x) < x - \dfrac{x^2}{2} + \dfrac{x^3}{3}$。
\end{examplebox}

\textbf{泰勒视角：}$\ln(1+x) = x - \dfrac{x^2}{2} + \dfrac{x^3}{3} - \dfrac{x^4}{4} + \cdots$，第四项为负，故截断后有严格小于关系。

\textbf{严格证明：}设 $f(x) = x - \dfrac{x^2}{2} + \dfrac{x^3}{3} - \ln(1+x)$，
\[
f'(x) = 1 - x + x^2 - \frac{1}{1+x} = \frac{(1-x+x^2)(1+x) - 1}{1+x} = \frac{x^3}{1+x}
\]
当 $x > 0$ 时，$f'(x) > 0$，又 $f(0) = 0$，故 $f(x) > 0$，即 $\ln(1+x) < x - \dfrac{x^2}{2} + \dfrac{x^3}{3}$。\qed

\begin{examplebox}[【例题5】高考真题用泰勒视角（2019全国II卷理数）]
已知函数 $f(x) = \ln x - \dfrac{x+1}{x-1}$，证明 $f(x) < 0$。
\end{examplebox}

\textbf{泰勒视角分析：}$\ln x$ 在 $x = 1$ 附近展开：$\ln x = (x-1) - \dfrac{(x-1)^2}{2} + \cdots$

而 $\dfrac{x+1}{x-1} = 1 + \dfrac{2}{x-1}$，在 $x = 1$ 附近趋向无穷，所以 $\ln x - \dfrac{x+1}{x-1}$ 在 $x = 1$ 附近为负。当 $x$ 远离 $1$ 时，$\dfrac{x+1}{x-1}$ 趋向 $1$，而 $\ln x$ 增长缓慢，仍为负。

\textbf{严格证明：}$f(x) < 0$ 等价于 $\ln x < \dfrac{x+1}{x-1}$。

当 $x > 1$ 时：设 $g(x) = (x-1)\ln x - (x+1)$，$g'(x) = \ln x + \dfrac{x-1}{x} - 1 = \ln x - \dfrac{1}{x}$。

$g''(x) = \dfrac{1}{x} + \dfrac{1}{x^2} > 0$，故 $g'$ 递增。$g'(1) = -1 < 0$，但 $\lim_{x \to +\infty} g'(x) = +\infty$，所以 $g'$ 有唯一零点 $x_0$。

$g(1) = -2 < 0$，可以证明 $g(x) < 0$ 对所有 $x > 1$ 成立（略去详细讨论）。

当 $0 < x < 1$ 时类似讨论。综合得 $f(x) < 0$。\qed

\begin{examplebox}[【例题6】泰勒视角秒杀切线问题]
求曲线 $y = e^x$ 在 $x = 0$ 处的切线方程，并用此切线证明 $e^x \geq 1 + x$。
\end{examplebox}

\textbf{解：}$f(0) = 1$，$f'(0) = 1$，切线方程为 $y = x + 1$。

设 $h(x) = e^x - x - 1$，$h'(x) = e^x - 1$。当 $x > 0$ 时 $h'(x) > 0$；当 $x < 0$ 时 $h'(x) < 0$。故 $h(x)$ 在 $x = 0$ 处取极小值 $h(0) = 0$，因此 $h(x) \geq 0$，即 $e^x \geq 1 + x$。\qed

\begin{tipbox}[切线放缩法]
函数 $f(x)$ 在 $x = a$ 处的切线 $y = f(a) + f'(a)(x-a)$ 就是泰勒展开的前两项。当 $f''(x) > 0$（凸函数）时，切线在曲线下方，即 $f(x) \geq f(a) + f'(a)(x-a)$。
\end{tipbox}

\begin{examplebox}[【例题7】放缩证明综合]
证明：当 $x > 0$ 时，$x - \dfrac{x^2}{2} < \ln(1+x) < x$。
\end{examplebox}

\textbf{泰勒视角：}$\ln(1+x) = x - \dfrac{x^2}{2} + \dfrac{x^3}{3} - \cdots$，截断到第一项得上界，截断到第二项得下界。

\textbf{严格证明：}

右边：设 $f(x) = x - \ln(1+x)$，$f'(x) = 1 - \dfrac{1}{1+x} = \dfrac{x}{1+x} > 0$（$x > 0$），故 $f(x) > f(0) = 0$。

左边：设 $g(x) = \ln(1+x) - x + \dfrac{x^2}{2}$，$g'(x) = \dfrac{1}{1+x} - 1 + x = \dfrac{x^2}{1+x} > 0$（$x > 0$），故 $g(x) > g(0) = 0$。\qed

\begin{examplebox}[【例题8】高考真题泰勒视角（2021新高考I卷）]
已知函数 $f(x) = x(1 - \ln x)$，讨论 $f(x)$ 的单调性并求最大值。
\end{examplebox}

\textbf{解：}$f'(x) = (1 - \ln x) + x \cdot \left(-\dfrac{1}{x}\right) = 1 - \ln x - 1 = -\ln x$。

当 $0 < x < 1$ 时 $f'(x) > 0$，递增；当 $x > 1$ 时 $f'(x) < 0$，递减。

最大值 $f(1) = 1 \times (1 - 0) = 1$。

\textbf{泰勒视角补充：}$\ln x$ 在 $x = 1$ 附近 $\approx (x-1) - \dfrac{(x-1)^2}{2}$，故
\[
f(x) = x(1 - \ln x) \approx x\left(1 - (x-1) + \frac{(x-1)^2}{2}\right) = x\left(2 - x + \frac{(x-1)^2}{2}\right)
\]
在 $x = 1$ 附近，$f(x) \approx 1 + 0 - \dfrac{(x-1)^2}{2}$，可见 $x = 1$ 是极大值点，且极值为 $1$。这与严格计算一致。

\begin{examplebox}[【例题9】泰勒展开与不等式综合]
证明：对任意 $x \in \mathbb{R}$，$e^x + e^{-x} \geq 2 + x^2$。
\end{examplebox}

\textbf{泰勒视角：}$e^x + e^{-x} = 2 + x^2 + \dfrac{x^4}{12} + \cdots$，后面全是非负项。

\textbf{严格证明：}设 $f(x) = e^x + e^{-x} - 2 - x^2$，
\[
f'(x) = e^x - e^{-x} - 2x, \quad f''(x) = e^x + e^{-x} - 2
\]
由均值不等式 $e^x + e^{-x} \geq 2$，等号当且仅当 $x = 0$。故 $f''(x) \geq 0$，$f'(x)$ 单调递增，$f'(0) = 0$。

当 $x < 0$ 时 $f'(x) < 0$，当 $x > 0$ 时 $f'(x) > 0$。故 $f(x)$ 在 $x = 0$ 处取极小值 $f(0) = 0$，因此 $f(x) \geq 0$。\qed

\begin{examplebox}[【例题10】高考压轴题展开视角]
已知 $f(x) = e^x - ax$（$a > 0$），若 $f(x) \geq 1$ 对所有 $x$ 成立，求 $a$ 的范围。
\end{examplebox}

\textbf{泰勒视角：}$e^x \approx 1 + x + \dfrac{x^2}{2}$，故 $f(x) \approx 1 + (1-a)x + \dfrac{x^2}{2}$。要使 $f(x) \geq 1$，需要 $(1-a)x + \dfrac{x^2}{2} \geq 0$ 对所有 $x$ 成立。当 $x < 0$ 且 $|x|$ 很大时，$\dfrac{x^2}{2}$ 主导，无问题；但当 $x$ 接近 $0$ 时，线性项 $(1-a)x$ 的符号很关键。直觉上 $a = 1$ 是临界值。

\textbf{严格证明：}$f'(x) = e^x - a$，令 $f'(x) = 0$ 得 $x = \ln a$。

$f(\ln a) = a - a\ln a = a(1 - \ln a)$。要求 $f(\ln a) \geq 1$，即 $a(1 - \ln a) \geq 1$。

设 $g(a) = a(1 - \ln a)$，$g'(a) = 1 - \ln a - 1 = -\ln a$。当 $0 < a < 1$ 时 $g'(a) > 0$，当 $a > 1$ 时 $g'(a) < 0$。$g(1) = 1$。

故 $g(a) \leq 1$，等号当且仅当 $a = 1$。因此 $a(1 - \ln a) \geq 1$ 的唯一解为 $a = 1$。\qed

\begin{examplebox}[【例题11】放缩证明综合提高]
证明：当 $x > 0$ 时，$\left(1 + \dfrac{1}{x}\right)^x < e < \left(1 + \dfrac{1}{x}\right)^{x+1}$。
\end{examplebox}

\textbf{泰勒视角：}$\ln\left(1 + \dfrac{1}{x}\right) = \dfrac{1}{x} - \dfrac{1}{2x^2} + \cdots$，故 $x \ln\left(1 + \dfrac{1}{x}\right) \approx 1 - \dfrac{1}{2x} < 1$，即 $\left(1 + \dfrac{1}{x}\right)^x < e$。

\textbf{严格证明（左边）：}设 $f(x) = x\ln\left(1 + \dfrac{1}{x}\right) - 1$，令 $t = \dfrac{1}{x}$，则 $f = \dfrac{\ln(1+t)}{t} - 1$。

设 $h(t) = \ln(1+t) - t$，$h'(t) = \dfrac{1}{1+t} - 1 = \dfrac{-t}{1+t} < 0$（$t > 0$），$h(0) = 0$，故 $h(t) < 0$，即 $\ln(1+t) < t$，$\dfrac{\ln(1+t)}{t} < 1$。\qed

右边类似，设 $g(x) = (x+1)\ln\left(1 + \dfrac{1}{x}\right) - 1 > 0$，可证。\qed

\begin{examplebox}[【例题12】泰勒展开与数列极限]
利用 $e^x$ 的泰勒展开，证明 $\left(1 + \dfrac{1}{n}\right)^n < e$ 且 $\lim_{n \to \infty}\left(1 + \dfrac{1}{n}\right)^n = e$。
\end{examplebox}

\textbf{证明：}由二项式展开：
\[
\left(1 + \frac{1}{n}\right)^n = \sum_{k=0}^{n} \binom{n}{k} \frac{1}{n^k} = \sum_{k=0}^{n} \frac{1}{k!} \cdot \frac{n(n-1)\cdots(n-k+1)}{n^k}
\]
每一项 $\dfrac{n(n-1)\cdots(n-k+1)}{n^k} < 1$，故 $\left(1 + \dfrac{1}{n}\right)^n < \displaystyle\sum_{k=0}^{n} \dfrac{1}{k!} < e$。

另一方面，对固定 $m \leq n$：
\[
\left(1 + \frac{1}{n}\right)^n \geq \sum_{k=0}^{m} \frac{1}{k!} \cdot \frac{n(n-1)\cdots(n-k+1)}{n^k}
\]
令 $n \to \infty$，右边趋于 $\displaystyle\sum_{k=0}^{m} \dfrac{1}{k!}$。再令 $m \to \infty$，得极限 $\geq e$。

综上，极限恰为 $e$。\qed

% ============ 第三部分：应用题 ============
\section*{\textcolor{apppurple}{三、应用题}}

\begin{appbox}[【应用1】工程中的线性近似]
某工厂生产零件，合格品的直径服从正态分布 $N(10, 0.04)$（单位 mm）。若公差要求为 $10 \pm 0.5$ mm，利用局部线性化的思想估算直径为 $10.1$ mm 时的概率密度相对于 $10$ mm 处的变化率。
\end{appbox}

\textbf{解：}$f(x) = \dfrac{1}{0.2\sqrt{2\pi}} e^{-\frac{(x-10)^2}{0.08}}$，$f'(x) = f(x) \cdot \left(-\dfrac{x-10}{0.04}\right)$。

$f'(10) = 0$（均值处变化率为零，这是正态分布的对称性决定的）。

更有趣的是看 $x = 10.1$ 处：$f'(10.1) = f(10.1) \times \left(-\dfrac{0.1}{0.04}\right) = -2.5 \times f(10.1)$。

局部线性化：$f(10.1 + \Delta x) \approx f(10.1) - 2.5 \times f(10.1) \times \Delta x$，说明在 $10.1$ mm 处，直径每增加 $0.1$ mm，概率密度约下降 $25\%$。

\begin{appbox}[【应用2】经济中的边际分析]
某企业生产 $x$ 件产品的总成本为 $C(x) = 0.01x^3 - 0.6x^2 + 15x + 100$（万元）。利用导数的局部线性化思想，估算从第 $20$ 件增加到第 $21$ 件时的边际成本，并与实际成本增量比较。
\end{appbox}

\textbf{解：}$C'(x) = 0.03x^2 - 1.2x + 15$。

边际成本 $C'(20) = 0.03 \times 400 - 1.2 \times 20 + 15 = 12 - 24 + 15 = 3$（万元）。

实际增量 $C(21) - C(20)$：
\begin{align*}
C(20) &= 0.01 \times 8000 - 0.6 \times 400 + 300 + 100 = 80 - 240 + 300 + 100 = 240 \\
C(21) &= 0.01 \times 9261 - 0.6 \times 441 + 315 + 100 = 92.61 - 264.6 + 315 + 100 = 243.01
\end{align*}
实际增量 $= 3.01$ 万元，与边际成本 $3$ 万元非常接近。这就是"局部线性化"在经济学中的应用——用导数近似增量。

\begin{appbox}[【应用3】物理中的小角度近似]
单摆的运动方程为 $\dfrac{d^2\theta}{dt^2} + \dfrac{g}{L}\sin\theta = 0$。利用泰勒展开 $\sin\theta \approx \theta$（小角度近似），求小振幅单摆的周期公式，并讨论近似的精度。
\end{appbox}

\textbf{解：}当 $\theta$ 很小时，$\sin\theta \approx \theta - \dfrac{\theta^3}{6}$。

零阶近似 $\sin\theta \approx \theta$：方程变为 $\dfrac{d^2\theta}{dt^2} + \dfrac{g}{L}\theta = 0$，周期 $T_0 = 2\pi\sqrt{\dfrac{L}{g}}$。

一阶修正：保留 $\theta^3$ 项，用摄动法可得 $T \approx T_0\left(1 + \dfrac{\theta_0^2}{16}\right)$，其中 $\theta_0$ 为振幅。

当 $\theta_0 = 15°\approx 0.262$ rad 时，修正量约 $\dfrac{0.262^2}{16} \approx 0.43\%$，误差很小。当 $\theta_0 = 30°$ 时，修正量约 $1.7\%$，仍然可接受。这就是泰勒展开在物理近似中的威力。

% ============ 第四部分：课堂练习 ============
\section*{\textcolor{tipgreen}{四、课堂练习}}

\begin{enumerate}[leftmargin=2em]
  \item 利用泰勒展开的前四项，计算 $\sin 0.5$ 的近似值（精确到小数点后四位）。

  \item 证明：当 $x > 0$ 时，$e^x > 1 + x + \dfrac{x^2}{2} + \dfrac{x^3}{6}$。

  \item 证明：当 $0 < x < 1$ 时，$\ln(1+x) > x - \dfrac{x^2}{2}$。

  \item 已知 $f(x) = e^x - \dfrac{1}{2}x^2 - x - 1$，证明 $f(x) \geq 0$。

  \item 利用泰勒展开证明：$\displaystyle\lim_{x \to 0} \frac{e^x - 1 - x}{x^2} = \dfrac{1}{2}$。

  \item 证明：当 $x > 0$ 时，$\tan x > x + \dfrac{x^3}{3}$。

  \item 已知 $f(x) = \ln(1+x) - x + ax^2$，若 $f(x) \leq 0$ 对所有 $x > 0$ 成立，求 $a$ 的最大值。

  \item 证明：$\displaystyle\sum_{k=1}^{n} \frac{1}{k} > \ln(n+1)$（提示：利用 $\ln(1+x) < x$）。
\end{enumerate}

\vspace{20pt}

\begin{knowledgebox}[参考答案]
\begin{enumerate}[leftmargin=2em]
  \item $\sin x \approx x - \dfrac{x^3}{6} + \dfrac{x^5}{120} - \dfrac{x^7}{5040}$，代入 $x = 0.5$：
  \[
  \sin 0.5 \approx 0.5 - \frac{0.125}{6} + \frac{0.03125}{120} - \frac{0.0078125}{5040} = 0.5 - 0.02083 + 0.00026 - 0.00000 = \boxed{0.4794}
  \]
  精确值 $\sin 0.5 \approx 0.47943$。

  \item 设 $f(x) = e^x - 1 - x - \dfrac{x^2}{2} - \dfrac{x^3}{6}$，$f'(x) = e^x - 1 - x - \dfrac{x^2}{2}$。由例题3知 $f'(x) > 0$（$x > 0$），又 $f(0) = 0$，故 $f(x) > 0$。\qed

  \item 设 $f(x) = \ln(1+x) - x + \dfrac{x^2}{2}$，$f'(x) = \dfrac{1}{1+x} - 1 + x = \dfrac{x^2}{1+x} > 0$（$x > 0$），$f(0) = 0$，故 $f(x) > 0$。\qed

  \item $f'(x) = e^x - x - 1$，$f''(x) = e^x - 1$。当 $x > 0$ 时 $f''(x) > 0$，$f'$ 递增，$f'(0) = 0$，故 $f'(x) > 0$。$f$ 递增，$f(0) = 0$，故 $f(x) \geq 0$。同理 $x < 0$ 时也成立。\qed

  \item $e^x = 1 + x + \dfrac{x^2}{2} + o(x^2)$，故 $e^x - 1 - x = \dfrac{x^2}{2} + o(x^2)$，$\dfrac{e^x - 1 - x}{x^2} = \dfrac{1}{2} + \dfrac{o(x^2)}{x^2} \to \boxed{\dfrac{1}{2}}$

  \item 设 $f(x) = \tan x - x - \dfrac{x^3}{3}$，$f'(x) = \sec^2 x - 1 - x^2 = \tan^2 x - x^2$。由 $\tan x > x$（$x > 0$），$f'(x) > 0$，$f(0) = 0$，故 $f(x) > 0$。\qed

  \item $f'(x) = \dfrac{1}{1+x} - 1 + 2ax = \dfrac{-x + 2ax(1+x)}{1+x} = \dfrac{x(2a + 2ax - 1)}{1+x}$。泰勒展开 $\ln(1+x) \approx x - \dfrac{x^2}{2}$，故 $f(x) \approx \left(a - \dfrac{1}{2}\right)x^2$。要使 $f(x) \leq 0$，需 $a \leq \dfrac{1}{2}$。严格验证 $a = \dfrac{1}{2}$ 时 $f(x) \leq 0$ 成立。$\boxed{a_{\max} = \dfrac{1}{2}}$

  \item 由 $\ln(1+x) < x$（$x > 0$），取 $x = \dfrac{1}{k}$：$\ln\left(1 + \dfrac{1}{k}\right) < \dfrac{1}{k}$，即 $\ln(k+1) - \ln k < \dfrac{1}{k}$。累加 $k = 1, 2, \ldots, n$：$\ln(n+1) < \displaystyle\sum_{k=1}^{n} \dfrac{1}{k}$。\qed
\end{enumerate}
\end{knowledgebox}

\vfill
\begin{center}
\small\textcolor{gray}{本讲义仅供教学参考，例题可根据学生水平调整难度。}
\end{center}

\end{document}
